% =========================================================
% Monotone Functions and Continuity — Notes
% Chapter: continuity
% Volume III · Analysis
% =========================================================

\subsection{Monotone Functions and Continuity}
\label{sec:monotone-functions}

\begin{tcolorbox}[title={Toolkit --- Monotone Functions}]
\begin{itemize}
  \item Every monotone function has $\operatorname{LeftLimit}$ and
    $\operatorname{RightLimit}$ at every interior point.
  \item $\operatorname{ContinuousAt}(f,c)\iff f(c^-)=f(c)=f(c^+)
    \iff j_f(c)=0$.
  \item All discontinuities of a monotone function are of first kind
    --- jump only, no oscillation.
  \item At most countably many discontinuities: disjoint jump
    intervals inject into $\mathbb{Q}$.
  \item Strictly monotone and continuous on $I$ $\Rightarrow$
    $f^{-1}$ exists, is strictly monotone, and is continuous.
  \item $\operatorname{ContinuousOn}(f,[a,b])\wedge\operatorname{Injective}(f)
    \iff f$ strictly monotone.
\end{itemize}
\end{tcolorbox}

\begin{theorem}[One-Sided Limits of Monotone Functions]
\label{thm:monotone-one-sided-limits}
Let $f$ be increasing on interval $I$, $c$ interior to $I$. Then:
\begin{enumerate}[label=(\roman*)]
  \item $\operatorname{LeftLimit}(f,c,\sup\{f(x):x<c\})$.
  \item $\operatorname{RightLimit}(f,c,\inf\{f(x):x>c\})$.
\end{enumerate}

\smallskip\noindent
\hyperref[prf:monotone-one-sided-limits]{\textit{Go to proof.}}
\end{theorem}

\begin{remark*}[Definition predicate reading]
\[
  f(c^-) \;:=\; \sup\{f(x):x\in I,\,x<c\},
  \qquad
  f(c^+) \;:=\; \inf\{f(x):x\in I,\,x>c\}.
\]
\end{remark*}

\begin{remark*}[Interpretation]
Monotonicity forces both one-sided limits to exist as finite values.
\end{remark*}

\begin{corollary}[Continuity Criterion for Monotone Functions]
\label{cor:monotone-continuity-criterion}
For $f$ increasing on $I$, $c$ interior:
$\operatorname{ContinuousAt}(f,c)\iff f(c^-)=f(c)=f(c^+)$.

\smallskip\noindent
\hyperref[prf:monotone-continuity-criterion]{\textit{Go to proof.}}
\end{corollary}

\begin{tcolorbox}[title={Definition --- Jump of a Function}]
\begin{definition}[Jump of a Function]
\label{def:jump-of-function}
For $f$ increasing on $I$, $c\in I$:
$j_f(c):=f(c^+)-f(c^-)$.
$\operatorname{ContinuousAt}(f,c)\iff j_f(c)=0$.
\end{definition}
\end{tcolorbox}

\begin{proposition}[Discontinuities of a Monotone Function Are of First Kind]
\label{prop:monotone-discontinuities-first-kind}
$\operatorname{MonotoneFunction}(f,E)\wedge\operatorname{DiscontinuousAt}(f,a)
\Rightarrow$ both $f(a^-)$, $f(a^+)$ exist finitely.

\smallskip\noindent
\hyperref[prf:monotone-discontinuities-first-kind]{\textit{Go to proof.}}
\end{proposition}

\begin{corollary}[Jump Intervals Are Disjoint]
\label{cor:jump-intervals-for-monotone-discontinuities}
For $f$ nondecreasing, distinct discontinuities $a\neq b$:
$(f(a^-),f(a^+))\cap(f(b^-),f(b^+))=\emptyset$.

\smallskip\noindent
\hyperref[prf:jump-intervals-for-monotone-discontinuities]{\textit{Go to proof.}}
\end{corollary}

\begin{theorem}[Discontinuities of a Monotone Function Are Countable]
\label{thm:monotone-discontinuities-countable}
$\operatorname{MonotoneFunction}(f,I)
\Rightarrow|\{c\in I:\operatorname{DiscontinuousAt}(f,c)\}|\leq\aleph_0$.

\smallskip\noindent
\hyperref[prf:monotone-discontinuities-countable]{\textit{Go to proof.}}
\end{theorem}

\begin{remark*}[Interpretation]
At each discontinuity $c$ the jump interval $(f(c^-),f(c^+))$ is
nonempty, open, and contains a rational. These intervals are pairwise
disjoint, giving an injection $D_f\hookrightarrow\mathbb{Q}$.
\end{remark*}

\begin{corollary}[Monotone Function Is Continuous on a Dense Set]
\label{cor:monotone-continuous-on-dense-set}
$\operatorname{MonotoneFunction}(f,[a,b])
\Rightarrow\{x:\operatorname{ContinuousAt}(f,x)\}$ is dense in $[a,b]$.

\smallskip\noindent
\hyperref[prf:monotone-continuous-on-dense-set]{\textit{Go to proof.}}
\end{corollary}

\begin{proposition}[Continuous Injective Is Strictly Monotone]
\label{prop:continuous-injective-iff-strictly-monotone}
$\operatorname{ContinuousOn}(f,[a,b])$:
$\operatorname{Injective}(f)\iff f$ strictly monotone.

\smallskip\noindent
\hyperref[prf:continuous-injective-is-strictly-monotone]{\textit{Go to proof.}}
\end{proposition}

\begin{tcolorbox}[colback=thmbox, colframe=thmborder, arc=2pt,
  left=6pt, right=6pt, top=4pt, bottom=4pt,
  title={\small\textbf{Theorem (Continuous Inverse Theorem)}},
  fonttitle=\small\bfseries]
\begin{theorem}[Continuous Inverse Theorem]
\label{thm:continuous-inverse-theorem}
$I$ an interval, $f$ strictly monotone and $\operatorname{ContinuousOn}(f,I)$
$\Rightarrow f^{-1}$ exists, is strictly monotone, and
$\operatorname{ContinuousOn}(f^{-1},f(I))$.

\smallskip\noindent
\hyperref[prf:inverse-strictly-monotone]{\textit{Go to proof.}}
\end{theorem}
\end{tcolorbox}


\begin{remark*}[Interpretation]
Continuity on $I$ forces $f(I)$ to be an interval. A jump
discontinuity of $f^{-1}$ would create a gap in $f(I)$ ---
contradiction. Hence $f^{-1}$ is continuous.
\end{remark*}
